\documentclass[12pt,a4paper]{ctexart}

% ==== Packages ====
\usepackage{amsmath,amssymb}
\usepackage[margin=2.5cm]{geometry}
\usepackage[hidelinks]{hyperref}
\usepackage{longtable}
\usepackage{booktabs}
\usepackage{array}
\usepackage{ragged2e}
\usepackage{fancyvrb}
\usepackage{fvextra}
\usepackage{xcolor}
\usepackage{enumitem}
\usepackage{graphicx}

% ==== Code block style ====
\DefineVerbatimEnvironment{shcode}{Verbatim}
  {fontsize=\small,frame=single,framesep=4pt,rulecolor=\color{gray!50},breaklines,breakanywhere}
\DefineVerbatimEnvironment{jsoncode}{Verbatim}
  {fontsize=\footnotesize,frame=single,framesep=4pt,rulecolor=\color{gray!50},breaklines,breakanywhere}

% ==== Metadata ====
\title{\textbf{Claude Code 完整配置指南\\[4pt]
\large DeepSeek 版}}
\author{}
\date{2026年7月}

\begin{document}
\maketitle

\begin{abstract}
\noindent 适合\textbf{零编程基础}用户，覆盖从 Node.js 安装、Git 配置、Claude Code 本体安装，到 DeepSeek API 连接、自定义 Skill 导入的全流程。全程约 20--30 分钟，Windows 和 Mac 用户均可轻松上手。特别补充了 Windows PowerShell 权限问题的解决方案，以及将他人分享的 Skill 能力包导入 Claude Code 的完整步骤。
\end{abstract}

\newpage

\tableofcontents

\newpage

\section{整体流程预览}

整个安装分为 6 步：

\begin{enumerate}[nosep]
    \item \textbf{安装 Node.js}（电脑的翻译器）
    \item \textbf{安装 Git}（文件管理器）
    \item \textbf{安装 Claude Code 本体}（AI 大脑）
    \item \textbf{安装编辑器}（VS Code 或 Trae）
    \item \textbf{配置 API 连接}（让 AI 连上 DeepSeek 服务器）
    \item \textbf{导入 Skill}（安装别人分享的专属 AI 能力包）
\end{enumerate}

\section{第一步：安装 Node.js}

Claude Code 需要 Node.js 作为运行环境，可以把它理解成一个``翻译器''。

\subsection{Windows}

\begin{enumerate}[nosep]
    \item 打开浏览器，访问 \url{https://nodejs.org/en/download/}
    \item 点击 \textbf{Windows Installer (.msi)} 下载
    \item 双击下载的文件，一路点 \textbf{Next}（所有默认选项都对）
    \item 验证安装：按 \texttt{Win} 键，输入 \textbf{PowerShell}，回车打开
\end{enumerate}

\begin{shcode}
node -v
\end{shcode}

看到类似 \texttt{v20.10.0} 就成功了。

\begin{shcode}
npm -v
\end{shcode}

\textbf{如果 npm -v 报红字错误（常见于首次使用 PowerShell）}，先输入：

\begin{shcode}
Set-ExecutionPolicy RemoteSigned
\end{shcode}

会提示输入 \texttt{Y} 确认，输完后再试 \texttt{npm -v} 看到版本号就正常了。

\subsection{Mac}

\begin{enumerate}[nosep]
    \item 访问 \url{https://nodejs.org/en/download/}
    \item 点击 \textbf{macOS Installer (.pkg)} 下载
    \item 双击安装，按提示操作
    \item 打开终端（在启动台搜索``终端''），验证：
\end{enumerate}

\begin{shcode}
node -v
npm -v
\end{shcode}

\section{第二步：安装 Git}

Git 是版本管理工具，Claude Code 需要它来管理文件。

\subsection{Windows}

\begin{enumerate}[nosep]
    \item 访问 \url{https://git-scm.com/downloads}
    \item 下载 Windows 版本，双击安装
    \item 大部分选项保持默认，重点注意：当看到 \textbf{``Adjusting your PATH environment''} 页面时，选择中间那个选项：\textbf{Git from the command line and also from 3rd-party software}
    \item 其他页面直接点 Next
\end{enumerate}

\begin{shcode}
git --version
\end{shcode}

看到版本号即成功。

\subsection{Mac}

Mac 通常自带 Git。打开终端输入：

\begin{shcode}
git --version
\end{shcode}

如果没装，系统会自动弹窗提示安装 Xcode Command Line Tools，点安装即可。

\section{第三步：安装 Claude Code 本体}

打开 \textbf{PowerShell}（Windows）或\textbf{终端}（Mac），输入：

\begin{shcode}
npm install -g @anthropic-ai/claude-code --registry=https://registry.npmmirror.com
\end{shcode}

\texttt{--registry=https://registry.npmmirror.com} 是使用国内镜像，下载更快。

等待几分钟，看到文字停止滚动且没有红色 \texttt{error} 字样，即安装成功。

\textbf{注意：}不要在命令前加 \texttt{sudo}，否则可能导致权限问题。

\section{第四步：安装编辑器}

编辑器是 Claude Code 的工作台，通过它和 AI 对话。二选一即可。

\subsection{选项一：VS Code（推荐）}

\begin{itemize}[nosep]
    \item 下载地址：\url{https://code.visualstudio.com/}
    \item Windows：下载后双击安装，建议勾选``添加到 PATH''和``添加右键菜单''
    \item Mac：下载后把 \texttt{Visual Studio Code.app} 拖到应用程序文件夹
\end{itemize}

\subsection{选项二：Trae IDE}

\begin{itemize}[nosep]
    \item 下载地址：\url{https://www.trae.cn/}
    \item 字节跳动旗下产品，基于 VS Code，界面对中文用户更友好
    \item 下载后双击安装即可
\end{itemize}

\section{第五步：配置 API 连接}

这是最关键的一步。

\subsection{安装 Claude Code 插件}

\begin{enumerate}[nosep]
    \item 打开 VS Code 或 Trae
    \item 点击侧边栏的\textbf{扩展图标}（四个方块图案）
    \item 搜索：\textbf{Claude Code for VS Code}
    \item 点击安装
\end{enumerate}

\subsection{获取 DeepSeek API Key}

\begin{enumerate}[nosep]
    \item 打开浏览器，访问 \url{https://platform.deepseek.com/}
    \item 注册 / 登录账号
    \item 进入 \textbf{API Keys} 页面，点击``创建新的 API Key''
    \item 复制保存这串 Key（格式类似 \texttt{sk-xxxxxxxxxxxxxxxx}）
\end{enumerate}

\textbf{重要：}API Key 相当于密码，请妥善保管，不要分享给他人或发到网上。

\subsection{写入配置文件}

\begin{enumerate}[nosep]
    \item 在 VS Code / Trae 中，按快捷键：
    \begin{itemize}[nosep]
        \item Windows：\texttt{Ctrl + Shift + P}
        \item Mac：\texttt{Cmd + Shift + P}
    \end{itemize}
    \item 输入：\textbf{Preferences: Open User Settings (JSON)}
    \item 回车，会打开一个 JSON 配置文件
    \item 将以下内容粘贴进去（如果文件已有其他内容，合并时注意 JSON 格式）：
\end{enumerate}

\begin{jsoncode}
{
  "env": {
    "ANTHROPIC_BASE_URL": "https://api.deepseek.com/anthropic",
    "ANTHROPIC_AUTH_TOKEN": "你的DeepSeek-API-Key粘贴在这里",
    "ANTHROPIC_MODEL": "DeepSeek-v4-pro[1m]",
    "ANTHROPIC_DEFAULT_HAIKU_MODEL": "DeepSeek-v4-flash",
    "ANTHROPIC_DEFAULT_SONNET_MODEL": "DeepSeek-v4-pro[1m]",
    "ANTHROPIC_DEFAULT_OPUS_MODEL": "DeepSeek-v4-pro[1m]",
    "ANTHROPIC_REASONING_MODEL": "DeepSeek-v4-pro[1m]",
    "ENABLE_TOOL_SEARCH": "true"
  },
  "includeCoAuthoredBy": false,
  "enabledPlugins": {
    "claude-hud@claude-hud": true,
    "skill-creator@claude-plugins-official": true
  },
  "extraKnownMarketplaces": {
    "claude-hud": {
      "source": {
        "source": "github",
        "repo": "jarrodwatts/claude-hud"
      }
    }
  },
  "effortLevel": "high"
}
\end{jsoncode}

\textbf{关键操作：}找到 \texttt{"你的DeepSeek-API-Key粘贴在这里"} 这一行，把中文部分替换成 DeepSeek API Key。\textbf{其他地方不要动。}

替换后类似：\texttt{"sk-xxxxxxxxxxxxxxxxxxxxxxxxxxxxxxxx"}

\subsection{保存并重启}

\begin{itemize}[nosep]
    \item 按 \texttt{Ctrl + S}（Windows）或 \texttt{Cmd + S}（Mac）保存
    \item \textbf{完全关闭} VS Code / Trae，再重新打开
    \item 重启后配置才能生效
\end{itemize}

\subsection{测试}

\begin{enumerate}[nosep]
    \item 打开 VS Code / Trae，随便打开一个文件夹
    \item 点击 Claude Code 图标（通常在侧边栏或右上角）
    \item 输入 \texttt{你好}
    \item AI 回复你了 $\to$ 配置成功！
\end{enumerate}

\section{第六步：导入 Skill（AI 能力包）}

Skill 是别人配置好的专属 AI 能力，导入后 Claude Code 会自动识别并加载。

\subsection{获取 Skill}

本仓库的技能直接来自 GitHub，不需要压缩包：

\begin{shcode}
git clone https://github.com/liuhao-hn/claude-multi-agent-workflows.git ~/claude-multi-agent-workflows
\end{shcode}

\subsection{安装}

\textbf{Mac / Linux——打开终端，输入：}

\begin{shcode}
mkdir -p ~/.claude/skills
cp -r ~/claude-multi-agent-workflows/skills/* ~/.claude/skills/
\end{shcode}

\textbf{Windows——打开 PowerShell，输入：}

\begin{shcode}
git clone https://github.com/liuhao-hn/claude-multi-agent-workflows.git "$env:USERPROFILE\claude-multi-agent-workflows"
New-Item -ItemType Directory -Force -Path "$env:USERPROFILE\.claude\skills"
Copy-Item "$env:USERPROFILE\claude-multi-agent-workflows\skills\*" "$env:USERPROFILE\.claude\skills\" -Recurse
\end{shcode}

只想装其中几个，把 \texttt{skills/*} 换成具体的技能目录名即可。也可以用软链（\texttt{ln -s}）指向仓库，这样改一次仓库本地立即生效。

\subsection{验证}

安装完成后，确保目录结构如下：

\begin{shcode}
~/.claude/skills/
├── coder-critic-review-team/
│   └── SKILL.md
├── commander-executor/
│   └── SKILL.md
├── gaodun-essay-grader/
│   └── SKILL.md
├── md2pdf/
│   └── SKILL.md
├── resume-generator/
│   └── SKILL.md
└── workflow-to-skill/
    └── SKILL.md
\end{shcode}

重启 Claude Code 后，输入对话触发相应功能即可自动加载 Skill。

\subsection{本仓库包含的能力}

{
\centering
\small
\begin{tabular}{>{\RaggedRight\arraybackslash}p{5.2cm}>{\RaggedRight\arraybackslash}p{9.3cm}}
\toprule
\textbf{Skill} & \textbf{说明} \\
\midrule
\texttt{commander-executor} & 多 Agent 分工总控（指挥官 + 外部执行者，TASKS.md 黑板 + 规则总线） \\
\midrule
\texttt{resume-generator} & 校招简历生成（模板版）：事实库 → 竞争策略 → 质检 → 单页中文 PDF \\
\midrule
\texttt{gaodun-essay-grader} & 申论大作文五维度批改（40 分制评分体系） \\
\midrule
\texttt{md2pdf} & MD → TEX → 专业中文 PDF 排版（ctexart + booktabs） \\
\midrule
\texttt{coder-critic-review-team} & 多 Agent 代码审查协作工作流 \\
\midrule
\texttt{workflow-to-skill} & 工作流自动沉淀引擎（任务结束 → 生成 Skill → 更新 Memory） \\
\bottomrule
\end{tabular}
\par}

想要更多 Skill？官方与第三方插件可用 \texttt{/plugin marketplace add} 直接安装，无需手动拷贝，见下一节。

\subsection{通过 Marketplace 安装更多 Skill（可选）}

除了上述打包好的 Skill，还可以从 GitHub Marketplace 安装官方维护的插件和 Skill。在 Claude Code 对话中直接输入以下命令：

\subsubsection*{Minimax Skills（14 个开发/文档/媒体 Skill）}

\begin{shcode}
/plugin marketplace add minimax-skills git https://github.com/MiniMax-AI/skills.git
\end{shcode}

\subsubsection*{Claude Code 官方插件（skill-creator、superpowers 等）}

\begin{shcode}
/plugin marketplace add claude-plugins-official github anthropics/claude-plugins-official
\end{shcode}

\subsubsection*{Claude HUD（终端状态栏）}

\begin{shcode}
/plugin marketplace add claude-hud github jarrodwatts/claude-hud
\end{shcode}

\subsubsection*{Stata Skill（Stata 代码参考）}

\begin{shcode}
/plugin marketplace add stata-skill github dylantmoore/stata-skill
\end{shcode}

添加 Marketplace 后，Skill 会自动加载可用，无需额外安装步骤。如需启用/禁用特定插件，在 settings.json 的 \texttt{enabledPlugins} 中配置。

\textbf{已安装的 Marketplace 一览：}

{
\centering
\small
\begin{tabular}{>{\RaggedRight\arraybackslash}p{3.5cm}>{\RaggedRight\arraybackslash}p{4.5cm}>{\RaggedRight\arraybackslash}p{6.5cm}}
\toprule
\textbf{Marketplace} & \textbf{来源} & \textbf{提供内容} \\
\midrule
minimax-skills & GitHub: MiniMax-AI/skills & 14 个 Skill（文档/开发/媒体） \\
\midrule
stata-skill & GitHub: dylantmoore/stata-skill & 2 个 Stata Skill \\
\midrule
claude-plugins-official & GitHub: anthropics/claude-plugins-official & 36 个插件（skill-creator、code-review、security-guidance 等） \\
\midrule
claude-hud & GitHub: jarrodwatts/claude-hud & 终端状态栏 HUD \\
\bottomrule
\end{tabular}
\par}

\vspace{8pt}

\textbf{汇总：}共 4 个 Marketplace，提供 16 个 Skill（MiniMax 14 + Stata 2）、36 个官方代码插件，以及 1 个终端状态栏 HUD。以上均在线安装，不随本仓库发布。

\newpage

\section{常见问题}

\subsection{npm install 提示权限错误（EACCES）}

不要用 sudo。按以下步骤修复：

\textbf{Mac：}

\begin{shcode}
mkdir ~/.npm-global
npm config set prefix '~/.npm-global'
echo 'export PATH=~/.npm-global/bin:$PATH' >> ~/.zshrc
source ~/.zshrc
\end{shcode}

重新打开终端，再次尝试安装。

\subsection{npm -v 报错（Windows PowerShell）}

\begin{shcode}
Set-ExecutionPolicy RemoteSigned
\end{shcode}

输入 \texttt{Y} 确认，然后重新执行 \texttt{npm -v}。

\subsection{API 连接失败}

\begin{enumerate}[nosep]
    \item 确认 API Key 复制正确，没有多余空格和换行
    \item 确认 settings.json 中的 JSON 格式正确（注意逗号和引号不能错）
    \item 确认已\textbf{完全关闭并重新打开}编辑器
    \item 检查网络连接是否正常（DeepSeek API 需要稳定的网络）
\end{enumerate}

\subsection{Skill 没有生效}

\begin{enumerate}[nosep]
    \item 确认文件夹放在了正确的路径（\texttt{\textasciitilde/.claude/skills/} 下）
    \item 每个 Skill 文件夹里必须有一个 \texttt{SKILL.md} 文件
    \item 重启 Claude Code
\end{enumerate}

\section{安装完成检查清单}

{
\centering
\begin{tabular}{>{\centering\arraybackslash}p{0.6cm}>{\RaggedRight\arraybackslash}p{13.9cm}}
\toprule
& \textbf{检查项} \\
\midrule
$\square$ & Node.js 已安装（\texttt{node -v} 显示版本号） \\
\midrule
$\square$ & Git 已安装（\texttt{git --version} 显示版本号） \\
\midrule
$\square$ & Claude Code 已安装（\texttt{claude --version} 无报错） \\
\midrule
$\square$ & VS Code 或 Trae 已安装 \\
\midrule
$\square$ & Claude Code 插件已安装 \\
\midrule
$\square$ & DeepSeek API Key 已配置到 settings.json \\
\midrule
$\square$ & 启动后输入``你好''得到 AI 回复 \\
\midrule
$\square$ & Skill 已安装到 \texttt{\textasciitilde/.claude/skills/} 目录（每个技能含 \texttt{SKILL.md}）\\
\bottomrule
\end{tabular}
\par}

\vspace{12pt}

全部打勾？恭喜，配置完成！

\end{document}
